\chapter{Host Institution: Polytech Intl}
\label{ch:host}

\section{Presentation}

This project was carried out under the supervision of Polytech Intl, within its Master's Program in Data Sciences. In the arrangement covered by this thesis, Polytech Intl acts as both the academic institution awarding the degree and the hosting and supervising body for the internship itself, rather than the internship being hosted by a separate company -- the program is structured so that students carry out a substantial, real-world capstone project under academic supervision, working directly with an external client organization on a genuine business problem rather than a simulated or purely academic exercise.

\section{Academic Supervision}

The project was supervised academically by Mr.~Marouene Chaieb, who provided guidance throughout the internship on both the technical direction of the work and its presentation as a piece of academic and professional deliverable work. Academic supervision in this arrangement covered methodology (ensuring the project's engineering process could be defended and evaluated, not just its output), scope discipline (helping keep a six-month project bounded to what could genuinely be delivered and verified), and the overall framing of the final deliverable -- this thesis.

\section{Why This Arrangement}

Structuring the internship this way -- university as host, external organization as client -- lets the program combine two things that are often in tension in a standard internship: genuine external accountability (the African Development Bank Group is a real client with a real problem, not a hypothetical case study) and the academic rigor of a supervised, methodologically grounded final thesis (CRISP-DM discipline, a running decision log, comprehensive automated testing) that a purely commercial internship placement does not always require or make time for. Chapter~2 describes the client side of this arrangement in detail.
